\newpage
\section*{Термины и определения}
\addcontentsline{toc}{section}{Термины и определения}

\begin{enumerate}
\item \textbf{VirtualBox (Oracle VM VirtualBox)} --- программный продукт
виртуализации для операционных систем Windows, Linux, macOS и других, позволяющий запускать гостевые операционные
системы в изолированной среде на базе хост-системы.

\item \textbf{Виртуализация} --- технология создания программного
(виртуального) представления вычислительных ресурсов (процессора, памяти,
накопителей, сетевых интерфейсов), позволяющая запускать несколько
независимых операционных систем на одном физическом устройстве.

\item \textbf{Виртуальная машина (ВМ)} --- программная эмуляция
вычислительного устройства, предоставляющая гостевой операционной системе
доступ к выделенным ей виртуальным аппаратным ресурсам.

% \item \textbf{Гипервизор} --- программный или программно-аппаратный
% компонент, обеспечивающий одновременное выполнение нескольких
% операционных систем на одном физическом компьютере и управляющий
% распределением ресурсов между ними.

% \item \textbf{VirtIO SCSI} --- паравиртуализированный драйвер контроллера
% накопителей, обеспечивающий взаимодействие гостевой операционной системы с
% виртуальным диском через оптимизированный интерфейс гипервизора,
% применяемый, в частности, для обеспечения совместимости при эмуляции
% процессорной архитектуры, отличной от архитектуры хост-системы.

\item \textbf{SLURM (Simple Linux Utility for Resource Management)} ---
система управления заданиями и ресурсами вычислительного кластера с
открытым исходным кодом, обеспечивающая постановку заданий в очередь, их
планирование и распределение по вычислительным узлам.

\item \textbf{Управляющий узел (head node, control node)} --- узел
вычислительного кластера, на котором выполняются диспетчер SLURM и
средства управления кластером, отвечающие за приём, планирование и
распределение заданий по вычислительным узлам.

\item \textbf{Вычислительный узел (compute node)} --- узел вычислительного
кластера, предоставляющий ресурсы (процессорное время, память) для
непосредственного выполнения заданий, поставленных в очередь диспетчера.

\item \textbf{SSH (Secure Shell)} --- сетевой протокол прикладного уровня,
позволяющий производить удалённое управление операционной системой и
туннелирование TCP-соединений, в том числе с использованием авторизации по
криптографическим ключам без передачи пароля по сети.

\item \textbf{Munge} --- служба аутентификации, используемая SLURM для
проверки подлинности сообщений, которыми обмениваются узлы кластера, на
основе общего секретного ключа.

\item \textbf{Разрешение имён} --- процесс сопоставления символьных имён
узлов сети соответствующим им IP-адресам, выполняемый средствами DNS или
локального файла соответствий (\texttt{/etc/hosts}).

\item \textbf{Заявка (задание, job)} --- описание вычислительной задачи,
поставленной пользователем в очередь диспетчера кластера, включающее
запрошенные ресурсы и оценку времени выполнения.

\item \textbf{Очередь заданий диспетчера} --- упорядоченный список заявок,
принятых диспетчером SLURM и ожидающих выделения ресурсов для запуска на
вычислительных узлах.

\item \textbf{Нормальное распределение} --- непрерывное распределение
вероятностей случайной величины, полностью определяемое математическим
ожиданием и дисперсией, график плотности которого имеет симметричную
колоколообразную форму.

\item \textbf{Логнормальное распределение} --- распределение вероятностей
случайной величины, логарифм которой имеет нормальное распределение с
параметрами $\mu$ и $\sigma$.

\item \textbf{Среднее арифметическое} --- статистическая характеристика
выборки, равная сумме всех её значений, делённой на их количество.

\item \textbf{Стандартное отклонение} --- статистическая характеристика
разброса значений случайной величины относительно её среднего
арифметического, равная квадратному корню из дисперсии.

\end{enumerate}